\documentclass[../../main.tex]{subfiles}%
\graphicspath{{\subfix{./image/}}{\subfix{../../image/}}}

\begin{document}

\section{参考文献与注记}\label{section:ch04-10}

我们关于 Sobolev 函数的主要依据是 Gilbarg--Trudinger\cite{gilbarg1983} 第 7 章与 Federer--Ziemer\cite{federer1972}。其中许多计算也出现在\cite{evans2010} 中。

关于光滑函数的磨光与局部逼近，见\cite{gilbarg1983} 第 7.2、7.3 节。\cref{theorem:ch04-4-2}取自\cite{gilbarg1983} 第 7.6 节，\cref{theorem:ch04-4-3}基于\cite{gilbarg1983} 中编号为 7.25 的定理。% XREF p=205 "Theorem 4.2 is from [G-T, Section 7.6]" | p=205 "and Theorem 4.3 is based upon [G-T, Theorem 7.25]"
乘积法则与链式法则见\cite{gilbarg1983} 第 7.4 节。关于延拓，另见\cite{gilbarg1983} 第 7.12 节。各种 Sobolev 型不等式见\cite{gilbarg1983} 第 7.7 节。\cref{section:ch04-05}中的\cref{lemma:ch04-4-1}是\cite{gilbarg1983} 中编号为 7.16 的引理的变体。% XREF p=205 "Lemma 4.1 in Section 4.5 is a variant of [G-T, Lemma 7.16]" | p=205 "Lemma 4.1 in Section 4.5 is a variant of [G-T, Lemma 7.16]"
紧性断言见\cite{gilbarg1983} 第 7.10 节。

在容量的处理上我们遵循\cite{federer1972}（参看 Maz'ja\cite{mazja1985} 与 Ziemer\cite{ziemer1989}）。\cref{section:ch04-07}中的\crefrange{theorem:ch04-4-14}{theorem:ch04-4-17}取自\cite{federer1972}，\cref{section:ch04-08}中的全部结果同样取自\cite{federer1972}。% XREF p=206 "Theorems 4.14–4.17 in Section 4.7 are from [F-Z]" | p=206 "Theorems 4.14–4.17 in Section 4.7 are from [F-Z]" | p=206 "as are all the results in Section 4.8"

关于容量的更多信息可在综合性著作\cite{ziemer1989} 与\cite{mazja1985} 中找到。Maly--Ziemer\cite{maly1997} 提供了在椭圆型偏微分方程解的正则性问题上的应用。Maly--Swanson--Ziemer\cite{maly2003} 讨论了 Sobolev 函数的余面积公式，Figalli\cite{figalli2008} 对 Sobolev 空间中的 Morse--Sard 定理给出了一个相当简单的证明。

\end{document}
